% Results excerpt. Compile from liquid_jax alongside paper.tex.
% Main figures: overall scaling; both ambiguity decompositions; task-error attribution.
% Full capacity curves and sensitivity analyses are supplied in supplement.md.

\paragraph{Scaling exhibits distinct performance regimes.}
We first average the five seeds within each capacity configuration and compute
its validation-loss reduction relative to $L=1$. We then average the matched
reductions over the 45 capacity configurations per task. Figure~\ref{fig:scaling-results}
includes fixed uniform routing ($L=0$) and reveals three regimes: a delayed
benefit on CIFAR-10, an intermediate optimum around 8--16 delegators on SVHN,
and early saturation followed by deterioration on the regression tasks.
Increasing from one to 32 delegators reduces direct validation loss by
$4.32\%$ on CIFAR-10 and $4.77\%$ on SVHN, but increases it by $13.64\%$
on Bikes and $4.32\%$ on Energy. The pointwise 95\% seed-block bootstrap
intervals for these reductions are $[3.78,4.93]\%$, $[3.00,5.60]\%$,
$[-15.32,-12.43]\%$, and $[-5.09,-3.58]\%$, respectively.
Superiority to one learned delegator does not necessarily imply superiority
to uniform routing: CIFAR-10 accuracy at 32 delegators remains 0.31 percentage
points below its uniform baseline, whereas SVHN exceeds it by 3.60 points.

\begin{figure}[t]
\centering
\includegraphics[width=\textwidth]{scaling_analysis/figures/01_scaling.pdf}
\caption{Final recorded validation-loss reduction relative to one delegator,
computed as $100(E_1-E_L)/E_1$ within each configuration after averaging
seeds, and then averaged across configurations; $E$ is recorded validation loss. Positive values favor scaling; $L=0$ is uniform
routing. Shading gives pointwise 95\% seed-block bootstrap intervals conditional
on the fixed grid and validation split.}
\label{fig:scaling-results}
\end{figure}

\paragraph{Capacity moderates the entire scaling response.}
We examine all three capacity factors across the complete scaling curve.
For each $L\in\{2,4,8,16,32\}$, an orthogonal factorial decomposition partitions
variation in the matched direct-loss gains across the balanced capacity grid
into main effects and interactions. Summing squared contributions across these
five levels shows that predictor width is the strongest individual moderator
on SVHN (38.0\% of observed capacity variation), but not universally.
On CIFAR-10, predictor count and delegator width contribute 26.8\% and 27.4\%,
respectively, while the predictor-width main effect contributes only 0.7\%.
Delegator width is the strongest single factor on Energy (33.6\%).
Interactions account for 41--54\% of the variation across tasks, making a
single-factor explanation insufficient. This partition describes the observed
grid, including finite-seed variability; it is not a causal importance measure.

Narrow SVHN predictors gain $11.96\%$ at eight delegators, compared with
$1.09\%$ for wide predictors. Bikes exhibits the opposite width dependence:
its 32-delegator penalty decreases from $18.73\%$ at predictor width four to
$9.36\%$ at width 16. Energy's wider predictors reach their useful scaling
limit earlier, whereas CIFAR-10's predictor-width curves largely converge
at large $L$. Thus limited predictor capacity does not have a uniform effect
on the value of delegator scaling.

\paragraph{Both components diversify while individual performance worsens.}
The performance--ambiguity decompositions in Figure~\ref{fig:scaling-members}
show that scaling increases recorded ambiguity in both predictors and
delegators, alongside increasing individual loss. Predictor error and
ambiguity nearly cancel, particularly in regression. On Bikes, mean
reference-weighted individual MSE rises from 16.65 to 97.00 between one and
32 delegators, while ambiguity rises from 16.60 to 96.95. Delegators likewise
become more diverse while their individual losses against the fitted routing
reference rise. These results indicate increasing compensation among members,
not improving individual correctness. As suggested by the error--diversity
perspective~\citep{wood2023unified}, the relevant question is whether the
resulting compensation lowers the ensemble's task loss.

\begin{figure}[t]
\centering
\includegraphics[width=\textwidth]{scaling_analysis/figures/03_loss_accounting.pdf}
\caption{Changes from $L=1$ in predictor (top) and delegator (bottom)
performance--ambiguity terms. With $\Delta X=X_L-X_1$, the top-row curves are
$100\Delta P/P_1$, $100\Delta A/P_1$, and $100\Delta(P-A)/P_1$, where $P$
and $A$ are the reference-weighted individual predictor loss and ambiguity.
For the bottom row, replace $P,A$ with individual routing loss $C$ and ambiguity
$D$, using $C_1$ as the common denominator. Seeds are averaged before computing
these percentages within each configuration; percentages are then averaged
across configurations. Negative green values indicate a decrease in the
loss-minus-ambiguity residual, not a negative residual.}
\label{fig:scaling-members}
\end{figure}

\paragraph{Task-loss attribution distinguishes the failure modes.}
We use the directly recorded identity $E=O+H$, where $O$ is predictor-ensemble
loss under the fitted reference delegators and $H=E-O$ is the signed excess
loss of learned delegation. All changes are normalized by the corresponding
$L=1$ task loss. Figure~\ref{fig:scaling-attribution} separates the scaling
effect into these two contributions.

On Bikes, the predictor/reference term improves by $11.17\%$ of baseline
loss, but the delegation term worsens by $24.81\%$, yielding a net $13.64\%$
penalty. Energy shows the same pattern: a $6.61\%$ predictor/reference
improvement is outweighed by a $10.93\%$ delegation penalty. Thus the
scaling-induced regression deterioration lies in delegation relative to the
fitted reference, despite improving predictor-ensemble capability under that
reference. Falling training MSE alongside rising validation MSE is consistent
with increasingly poor generalization of learned routing.

SVHN's late downturn has a different source. Between eight and 32 delegators,
predictor/reference loss increases by $2.57\%$ of the one-delegator loss,
while improved delegation offsets $1.04\%$, leaving a $1.53\%$ deterioration.
On CIFAR-10, the initial four-delegator penalty is almost entirely
delegation-side; at 32 delegators, both terms improve, with delegation
providing $3.28\%$ of the total $4.32\%$ gain.

\begin{figure}[t]
\centering
\includegraphics[width=\textwidth]{scaling_analysis/figures/04_fault_attribution.pdf}
\caption{Contributions to direct task-loss reduction relative to $L=1$.
For each configuration, the component contributions are
$100(O_1-O_L)/E_1$ and $100(H_1-H_L)/E_1$, using the same baseline total
task loss $E_1$; these sum to $100(E_1-E_L)/E_1$. Seed means precede
normalization, and percentages are then averaged across configurations.
Positive contributions improve performance; negative contributions worsen it.
The fitted reference changes with ensemble size, so this is conditional
error attribution rather than a controlled causal intervention.}
\label{fig:scaling-attribution}
\end{figure}

These attributions also explain the capacity dependence. On SVHN, narrow
predictors obtain a $9.26\%$ delegation improvement at 32 delegators, compared
with $0.33\%$ for wide predictors. On Bikes, widening predictors reduces the
delegation penalty from $30.51\%$ to $19.87\%$. Widening delegators improves
the delegation contribution on CIFAR-10 but amplifies the penalty on Energy.
The benefit therefore depends on whether learned routing generalizes well
enough to exploit the available predictor ensemble.

\paragraph{Measurement scope.}
All results concern the recorded validation split. The all-$L$ performance
curves use the logged objective; direct CE/MSE evaluation confirms the
positive-$L$ conclusions and supplies the task-loss attribution. The member
ambiguity terms use probability clipping, so their stored residuals are not
substituted for direct task loss. Predictor member statistics use the better
of learned and oracle weights, whereas $O$ and $H$ retain the fitted oracle
and signed gap. The oracle is validation-selected and its capacity varies
with $L$; the attribution is conditional on this reference. Increasing $L$
also increases parameters and compute, rather than holding the budget fixed.
